\documentclass[sigconf,nonacm]{acmart}
%% NOTE that a single column version may be required for 
%% submission and peer review. This can be done by changing
%% the \doucmentclass[...]{acmart} in this template to 
%% \documentclass[manuscript,screen]{acmart}
%% 
%% To ensure 100% compatibility, please check the white list of
%% approved LaTeX packages to be used with the Master Article Template at
%% https://www.acm.org/publications/taps/whitelist-of-latex-packages 
%% before creating your document. The white list page provides 
%% information on how to submit additional LaTeX packages for 
%% review and adoption.
%% Fonts used in the template cannot be substituted; margin 
%% adjustments are not allowed.
%%
%%
%% \BibTeX command to typeset BibTeX logo in the docs
\AtBeginDocument{%
  \providecommand\BibTeX{{%
    \normalfont B\kern-0.5em{\scshape i\kern-0.25em b}\kern-0.8em\TeX}}}

%% Rights management information.  This information is sent to you
%% when you complete the rights form.  These commands have SAMPLE
%% values in them; it is your responsibility as an author to replace
%% the commands and values with those provided to you when you
%% complete the rights form.
\setcopyright{acmcopyright}
\copyrightyear{2022}
\acmYear{2022}
\acmDOI{XXXXXXX.XXXXXXX}

%% These commands are for a PROCEEDINGS abstract or paper.
\acmConference[SCORES'22]{Student Computing Research Symposium}{October 6,
2022}{Ljubljana, Slovenia}
%
%  Uncomment \acmBooktitle if th title of the proceedings is different
%  from ``Proceedings of ...''!
%
%\acmBooktitle{Woodstock '18: ACM Symposium on Neural Gaze Detection,
%  June 03--05, 2018, Woodstock, NY} 
\acmPrice{15.00}
\acmISBN{978-1-4503-XXXX-X/18/06}


%%
%% Submission ID.
%% Use this when submitting an article to a sponsored event. You'll
%% receive a unique submission ID from the organizers
%% of the event, and this ID should be used as the parameter to this command.
%%\acmSubmissionID{123-A56-BU3}

%%
%% For managing citations, it is recommended to use bibliography
%% files in BibTeX format.
%%
%% You can then either use BibTeX with the ACM-Reference-Format style,
%% or BibLaTeX with the acmnumeric or acmauthoryear sytles, that include
%% support for advanced citation of software artefact from the
%% biblatex-software package, also separately available on CTAN.
%%
%% Look at the sample-*-biblatex.tex files for templates showcasing
%% the biblatex styles.
%%

%%
%% The majority of ACM publications use numbered citations and
%% references.  The command \citestyle{authoryear} switches to the
%% "author year" style.
%%
%% If you are preparing content for an event
%% sponsored by ACM SIGGRAPH, you must use the "author year" style of
%% citations and references.
%% Uncommenting
%% the next command will enable that style.
%%\citestyle{acmauthoryear}

\newcommand{\scores}{\textsc{SCORES}}
\renewcommand{\keywordsname}{Ključne besede}
\renewcommand{\acksname}{Zahvala}

\usepackage[slovene]{babel}
\usepackage{verbatim}
\usepackage{booktabs}
\usepackage[linesnumbered, boxed, resetcount]{algorithm2e}
\SetAlgorithmName{Algoritem}{Algoritem}

%%
%% If you need any other LaTeX packages, add them here
%%

%%
%% end of the preamble, start of the body of the document source.
\begin{document}

%%
%% The "title" command has an optional parameter,
%% allowing the author to define a "short title" to be used in page headers.
\title{Klasifikacija neželjenih SMS sporočil}


\author{David Krajnc}
\email{david.krajnc4@student.um.si}
\affiliation{%
    \institution{Faculty of Electrical \\Engineering and Computer Science,\\
    University of Maribor\\Koroška cesta 46}
    \city{SI-2000 Maribor}
    \country{Slovenia}
}

\author{David Lipavec}
\email{david.lipavec@student.um.si}
\affiliation{%
    \institution{Faculty of Electrical \\Engineering and Computer Science,\\
    University of Maribor\\Koroška cesta 46}
    \city{SI-2000 Maribor}
    \country{Slovenia}
}

\author{Matija Pajenk}
\email{matija.pajenk@student.um.si}
\affiliation{%
    \institution{Faculty of Electrical \\Engineering and Computer Science,\\
    University of Maribor\\Koroška cesta 46}
    \city{SI-2000 Maribor}
    \country{Slovenia}
}

\author{Andraž Šošterič}
\email{andraz.sosteric@student.um.si}
\affiliation{%
    \institution{Faculty of Electrical \\Engineering and Computer Science,\\
    University of Maribor\\Koroška cesta 46}
    \city{SI-2000 Maribor}
    \country{Slovenia}
}

\author{Tadej Teršek}
\email{tadej.tersek@student.um.si}
\affiliation{%
    \institution{Faculty of Electrical \\Engineering and Computer Science,\\
    University of Maribor\\Koroška cesta 46}
    \city{SI-2000 Maribor}
    \country{Slovenia}
}


\begin{abstract}
    Povzetek
\end{abstract}


\keywords{SMS, vsiljiva sporočila, klasifikacija}


\maketitle

\section{Uvod}
    Mobilni telefoni so skozi leta postali del našega vsakdanjega življenja. Njihova
    uporaba iz leta v leto narašča, leta 2025 pa je mobilni telefon uporabljalo kar 
    5.28 milijard uporabnikov, kar je 8\% več kot leta 2024 \cite{HowManyPhones}. Kljub različnim aplikacijam
    za sporočanje, kot so Facebook, WhatsApp, Viber ipd., so SMS sporočila še vedno 
    priljubljena, saj omogočajo hitro in enostavno komunikacijo. V letu 2025 je bilo
    poslano okoli 25 milijard SMS sporočil dnevno \cite{HowManySMS}. Vendar pa je veliko teh sporočil 
    neželenih, saj jih pošiljajo različni ponudniki storitev in prevaranti, da bi nam 
    prodali svoje izdelke ali storitve. Ravno zaradi tega je pomembno, da imamo učinkovite 
    metode za klasifikacijo SMS sporočil, da lahko ločimo med želenimi in neželenimi 
    sporočili. Cilj tega članka je predstaviti štiri različne metode klasifikacije SMS 
    sporočil, ki temeljijo na različnih algoritmih strojnega učenja, kjer bomo primerjali
    njiihovo natančnost, hitrost in učinkovitost. Za učenje modela bomo uporabili javno 
    dostopen nabor podatkov SMS sporočil, ki vsebuje tako 28560 želenih kot 18832 neželenih 
    oz. vsiljivih sporočil \cite{DatasetSpamSMS}.
    V nadaljevanju članka bomo najprej predstavili sorodna dela na področju klasifikacije 
    SMS sporočil, nato pa bomo opisali našo lastno idejo in pristop k reševanju problema 
    klasifikacije SMS sporočil. Na koncu bomo predstavili rezultate naših eksperimentov in 
    zaključek.


\section{Sorodna dela}
    V zadnjem desetletju je tematika klasifikacije nezaželenih SMS sporočil pritegnila precej pozornosti raziskovalcev. Razlog je predvsem v tem, da se spam sporočila razlikujejo od klasične elektronske pošte \cite{lota2017systematic}. Gre predvsem za omejeno dolžino besedila in specifičen slog jezika, denimo uporabo kratic ali emojijev.

    Avtorji Xu in sodelavci so se v svoji raziskavi osredotočili na detekcijo brez uporabe vsebine sporočil (angl. non-content features). Namesto analize besedila so preučevali metapodatke in vedenjske vzorce, kar dokazuje, da je mogoče visoko natančnost doseči tudi z metodami, ki ne posegajo v zasebnost vsebine sporočil, temveč analizirajo le omrežne interakcije \cite{6133257}.
    
    Po drugi strani pa se večina raziskav opira na algoritme strojnega učenja, ki analizirajo neposredno vsebino. Sethi in sodelavci so izvedli primerjalno analizo različnih algoritmov, kjer so poudarili pomen predobdelave podatkov in uporabo tehnik naravnega jezika (NLP) za ekstrakcijo značilk \cite{8284445}. Njihovi rezultati potrjujejo, da so modeli, kot so naivni Bayesov klasifikator in podporni vektorji (SVM), izjemno učinkoviti pri binarni klasifikaciji sporočil.
    
    Glavno motivacijo za pričujočo seminarsko nalogo predstavlja delo Shirani-Mehra. Avtor je v svoji raziskavi sistematično preučil uporabo strojnega učenja za detekcijo spama, pri čemer je primerjal zmogljivost več modelov na standardnih naborih podatkov \cite{shirani2013sms}. Njegov pristop služi kot temelj za našo implementacijo. V tem delu se osredotočamo na rekreacijo njegovih eksperimentov z namenom preverjanja ponovljivosti rezultatov v sodobnem okolju.


\section{Lastna ideja}
    Lastna ideja

\bibliographystyle{ACM-Reference-Format}
\bibliography{bibliografija}

\end{document}
